\documentclass[11pt, letterpaper]{article}

\usepackage[proposal]{POL380}   % paper: [paper]
\usepackage{natbib} % For bibliography 

%%%%%%%%%%%%%%%% Title and Heading  %%%%%%%%%%%%%%%%%%%%%%%

\papertitle{Identifying the Causes of Protectionism using Double Machine Learning }
\paperauthor{Emily Harms}
\paperaffil{The College of New Jersey}
\paperdate{\today}

\begin{document}
\maketitle

%%%%%%%%%%%%%%%%  Topic Description  %%%%%%%%%%%%%%%%%%%%%%%

\section{Proposed Topic}

Why do some democracies protect domestic producers more than others? I aim to study whether factors such as 
exports, unemployment rates, GDP, inflation rates, and the number of political parties causally shape how protectionist a country's trade policy is. Using a panel of many countries spanning several decades, I will estimate the effect of each factor using double machine learning (DML), which lets flexible learners such as random forests and neural networks estimate causal parameters in the presence of complex confounding variables \citep{DML}.  

This project aims to answer the question \textbf{Which domestic factors cause countries to adopt more protectionist trade policies, and how robust are these effects?} by testing a variety of hypotheses regarding protectionism. \textcolor{red}{Write a few sentences on why each is expected and the alternate is not expected} 

\paragraph{H1 (Exports): } Exports being a higher percentage of GDP \textcolor{red}{increases/decreases -- which makes more sense?} protectionism. 

\paragraph{H2 (Umeployment):} Increased unemployment rates cause increased protectionism.

\paragraph{H3 (GDP per Capita):} Increases in GDP per capita causes increased protectionism

\paragraph{H4 (Total GDP in US\$):} A higher total GDP causes increased protectionism. 

\paragraph{H5 (Inflation):} Increased inflation rates cause increased protectionism.

\paragraph{H6 (Number of Political Parties):} More political parties causes increased protectionism. \\

\noindent For the sake of this analysis, the level of protectionism will be quantified using tariff rates \textcolor{red}{explain why this is a reasonable thing to do.} In other words, higher tariff rates mean higher levels of protectionism, and vice versa. 

%%%%%%%%%%%%%%%%  Prior Work  %%%%%%%%%%%%%%%%%%%%%%%


\section{Prior Work}

%%%%%%%%%%%%%%%%  Methods/Progress  %%%%%%%%%%%%%%%%%%%%%%%

\section{Proposed Methods and Current Progress} 

\noindent \textcolor{red}{Explain DML in detail, say how it works, which algorithms specifically, why those algorithms (likely random forest and neural networks). Make an argument for why DML instead of regression.}  

Data on tariff rates \citep{TariffData}, exports \citep{ExportData}, unemployment rates \citep{UnemploymentData}, GDP \citep{GDP_US_Data, GDPCapitaData}, and inflation \citep{InflationData} and the number of political parties \citep{NoPartiesData} have been obtained, and the required data reduction and cleaning is complete. A simple Python program \textcolor{red}{link to my GitHub} was written to pull the data needed from each dataset, and add the data to one large .csv file. A subset of the data is shown in table 1. \\

\begin{center}
\begin{table}[h]
\makebox[\textwidth]{
\begin{tabular}{c c c c c c c c c c}
\hline\hline
Country & Year & \% Exports & Unemp. & GDP/Capita & GDP (US\$) & Inflation & No. Parties & Tariff  \\ [0.5ex]
\hline
ALB & 1997 & 10.04 & 16.874 & 728.55 & 4.71e9 & 33.18 & 2.87 & 14.41 \\
ALB & 2000 & 19.53 & 19.023 & 1160.42 & 6.15e9 & 0.05 & 2.87 & 11.65 \\
... & ... & ... & ... & ... & ... & ... & ... &  \\
ZMB &  2020 & 46.72 & 6.032 & 951.64 & 2.34e10 & 15.73 & 2.78 & 4.82 \\ 
\hline
\end{tabular} }
\caption{Dataset that will be used for double machine learning. Country and Year are for reference \textcolor{red}{though it may be interesting to look at each country's level of protectionism over time}, year, \% exports, unemployment rate, GDP/capita, GDP in US\$, inflation rate, and number of political parties are independent variables, tariff rate is the dependent variable.}
\end{table}
\end{center}

%%%%%%% Preliminary Results %%%%%%%%%%

\begin{center}
\begin{table}[h]
\makebox[\textwidth]{
\begin{tabular}{c c c c c c c}
\hline\hline
Parameter & coef & std err & t & P > |t| & 2.5 \% & 97.5 \% \\ [0.5ex]
\hline
Exports & -0.0227 & 0.00962 & -2.363 & 0.0181 & -0.0416 & -0.00388 \\
Unemployment & 0.0279 & 0.0511 & 0.546 & 0.585 & -0.0722 & 0.128 \\ 
GDP/Capita & -3.898 & 8.559e-6 & -4.554 & 5.254e-6 & -5.483e-5 & -2.221e-5 \\
GDP (US\$) & -2.473e-13 & 0.00153 & -1.610 & 0.107 &  -5.483e-13 & 5.375e-14 \\ 
Inflation & -5.581e-4 & 0.00153 & -0.366 & 0.715 & -0.00355 & 0.00243 \\
Number of Parties & 0.0679 & 0.0725 & 0.936 & 0.350 & -0.0743 & 0.210 \\
\hline
\end{tabular} }
\caption{Preliminary results. A P-value less than 0.5 shows a statistically significant causal relationship.}
\end{table}
\end{center}


%%%%%%%%%%%%%%%%  Impacts of Protectionism  %%%%%%%%%%%%%%%%%%%%%%%

\section{Impacts of Protectionism} 

%%%%%%%%%%%%%%%% Notes %%%%%%%%%%%%%%%%%%%%%%%

\section{Notes}

\begin{itemize}
	\item Check DML Paper citation --> May be linking to ArXiv, find actual publication if so
	\item Is JeanReshef2017 peer reviewed? Double check.
	\item Check years for World Bank data and V-Dem data 
\end{itemize}


%%%%%%%%%%%%%%%% Bibliography  %%%%%%%%%%%%%%%%%%%%%%%

\bibliographystyle{apsa-leeper}
\bibliography{POL380_Proposal}

\end{document}